\chapter{Evaluation Methodology}

\section{Overview}

The evaluation of TextLens was conducted based on the actual execution pipeline of the implemented system rather than using a generic OCR benchmark. The application processes camera frames continuously and performs several stages of analysis, including motion stability detection, OCR execution, language detection, translation, tracking, and rendering. Since these stages occur asynchronously and operate on different frame intervals, the evaluation was structured to measure the performance of each stage individually as well as the overall user-perceived latency.

All timing measurements were collected using internal timestamps generated by the application during execution. These timestamps were recorded at key points in the pipeline, including camera initialization, frame acquisition, stable frame confirmation, OCR completion, translation completion, and overlay rendering. Using internal timestamps allowed accurate measurement of the delay experienced by the user while preserving the event-driven nature of the application.

The evaluation focused on five main aspects of the system:

\begin{itemize}
\item OCR detection and recognition accuracy
\item Language routing accuracy
\item Translation quality
\item System latency
\item Real-time throughput
\end{itemize}

Each of these components was measured using metrics appropriate to its function within the pipeline.

\section{OCR Accuracy Evaluation}

OCR performance was evaluated in two stages: text detection accuracy and text recognition accuracy.

Text detection accuracy was measured using precision and recall. Precision measures the proportion of detected text regions that correspond to actual text, while recall measures the proportion of real text regions that were successfully detected by the system.

\begin{equation}
Precision = \frac{TP}{TP + FP}
\end{equation}

\begin{equation}
Recall = \frac{TP}{TP + FN}
\end{equation}

Where:
\begin{itemize}
\item $TP$ represents correctly detected text regions
\item $FP$ represents incorrectly detected regions
\item $FN$ represents missed text regions
\end{itemize}

Text recognition accuracy was measured using Character Error Rate (CER), which measures the difference between the OCR output and the ground truth transcription.

\begin{equation}
CER = \frac{S + D + I}{N}
\end{equation}

Where:
\begin{itemize}
\item $S$ is the number of substituted characters
\item $D$ is the number of deleted characters
\item $I$ is the number of inserted characters
\item $N$ is the total number of characters in the reference text
\end{itemize}

Lower CER values indicate better recognition performance.

\section{Language Routing Evaluation}

After text recognition, the system performs language detection to determine which translation model should be used. Since different models are used for Latin-based languages and CJK languages, accurate language routing is essential to ensure correct translation.

Language routing accuracy was measured as the percentage of text segments that were correctly classified into their respective language groups.

\begin{equation}
Routing\ Accuracy = \frac{Correct\ Classifications}{Total\ Text\ Segments}
\end{equation}

During evaluation, OCR outputs containing multiple languages were manually labeled and compared against the routing decisions made by the language detection module.

\section{Translation Quality Evaluation}

Translation quality was evaluated using manual assessment rather than automated metrics such as BLEU or METEOR. This approach was chosen because the application processes short scene text segments, such as signs and labels, where semantic correctness and readability are more important than sentence-level n-gram overlap.

A sample set of translated text segments was manually reviewed and rated using a five-point scale:

\begin{itemize}
\item 5 – Perfect translation with correct meaning and grammar
\item 4 – Minor grammatical issues but meaning preserved
\item 3 – Understandable translation with noticeable errors
\item 2 – Significant translation errors affecting meaning
\item 1 – Incorrect or unusable translation
\end{itemize}

The final translation score was calculated as the average rating across all evaluated samples.

\section{Latency Measurement}

System latency represents the delay between the moment the camera observes a stable scene and the moment the translated overlay becomes visible on the screen.

Latency was measured using timestamps recorded at the following stages:

\begin{itemize}
\item Stable frame detection
\item OCR completion
\item Translation completion
\item Overlay rendering
\end{itemize}

The total latency was computed as:

\begin{equation}
Total\ Latency = OCR\ Time + Translation\ Time + Rendering\ Time
\end{equation}

The measurements were repeated multiple times across different scenes to obtain an average latency value.

\section{Throughput Evaluation}

Throughput was evaluated to determine whether the application could maintain real-time performance while running continuously on a mobile device.

Because OCR and translation are computationally expensive operations, they were executed only on stable frames rather than on every captured frame. Between OCR refresh cycles, the tracking module propagated translated text positions across new frames.

Three throughput metrics were measured:

\begin{itemize}
\item Camera frame processing rate
\item OCR processing rate
\item Overlay rendering frame rate
\end{itemize}

These metrics allowed the system to maintain a smooth visual output while limiting the frequency of expensive OCR computations.

\section{Test Environment}

All experiments were conducted using the deployed Android application running on a mobile device under typical indoor lighting conditions. The application processed live camera input and evaluated scenes containing printed signs, labels, and multi-language text.

Test scenes included Latin alphabet languages as well as Chinese, Japanese, and Korean text in order to verify the correct operation of the language routing mechanism.

Each test scenario was repeated multiple times to reduce variability caused by camera motion and environmental factors.